\documentclass[11pt]{article}
% Shared preamble for all ML-RESEARCH papers.
% Papers load it with: % Shared preamble for all ML-RESEARCH papers.
% Papers load it with: % Shared preamble for all ML-RESEARCH papers.
% Papers load it with: \input{../../../papers-common/preamble.tex}
\usepackage[utf8]{inputenc}
\usepackage[T1]{fontenc}
\usepackage{lmodern}
\usepackage{amsmath,amssymb,amsthm}
\usepackage{graphicx}
\usepackage{booktabs}
\usepackage{hyperref}
\usepackage[margin=1.1in]{geometry}

\newtheorem{theorem}{Theorem}
\newtheorem{proposition}{Proposition}
\newtheorem{definition}{Definition}
\newtheorem{remark}{Remark}

\newcommand{\R}{\mathbb{R}}
\newcommand{\E}{\mathbb{E}}
\newcommand{\Prob}{\mathbb{P}}
\newcommand{\ip}[2]{\langle #1, #2 \rangle}

\usepackage[utf8]{inputenc}
\usepackage[T1]{fontenc}
\usepackage{lmodern}
\usepackage{amsmath,amssymb,amsthm}
\usepackage{graphicx}
\usepackage{booktabs}
\usepackage{hyperref}
\usepackage[margin=1.1in]{geometry}

\newtheorem{theorem}{Theorem}
\newtheorem{proposition}{Proposition}
\newtheorem{definition}{Definition}
\newtheorem{remark}{Remark}

\newcommand{\R}{\mathbb{R}}
\newcommand{\E}{\mathbb{E}}
\newcommand{\Prob}{\mathbb{P}}
\newcommand{\ip}[2]{\langle #1, #2 \rangle}

\usepackage[utf8]{inputenc}
\usepackage[T1]{fontenc}
\usepackage{lmodern}
\usepackage{amsmath,amssymb,amsthm}
\usepackage{graphicx}
\usepackage{booktabs}
\usepackage{hyperref}
\usepackage[margin=1.1in]{geometry}

\newtheorem{theorem}{Theorem}
\newtheorem{proposition}{Proposition}
\newtheorem{definition}{Definition}
\newtheorem{remark}{Remark}

\newcommand{\R}{\mathbb{R}}
\newcommand{\E}{\mathbb{E}}
\newcommand{\Prob}{\mathbb{P}}
\newcommand{\ip}[2]{\langle #1, #2 \rangle}


\title{The Linear Discriminator as a Bernoulli Hypothesis Class}
\author{Yuval Lavie}
\date{\today}

\begin{document}
\maketitle

\begin{abstract}
We define the linear discriminator as a hypothesis class of conditional
Bernoulli distributions and show that maximum-likelihood estimation over this
class is exactly minimization of the logistic loss. This note is the
theoretical companion to the implementation in
\texttt{mlr.models.pure.linear\_discriminator}.
\end{abstract}

\section{The hypothesis class}

Let $\mathcal{X} = \R^d$ and $\mathcal{Y} = \{0, 1\}$. For a weight vector
$w \in \R^d$ and bias $b \in \R$, define the sigmoid
$\sigma(z) = (1 + e^{-z})^{-1}$ and the hypothesis
\begin{equation}
  h_{w,b}(x) \;=\; \sigma\!\left(\ip{w}{x} + b\right) \in (0, 1).
\end{equation}

\begin{definition}[Bernoulli hypothesis class]
The linear discriminator class is
$\mathcal{H} = \{\, h_{w,b} : w \in \R^d,\; b \in \R \,\}$, where each
hypothesis is read as a conditional distribution
$\Prob(Y = 1 \mid X = x) = h_{w,b}(x)$, i.e.\
$Y \mid X = x \sim \mathrm{Bernoulli}\!\left(h_{w,b}(x)\right)$.
\end{definition}

\section{Maximum likelihood is logistic-loss minimization}

\begin{proposition}
Given a sample $S = \{(x_i, y_i)\}_{i=1}^n$, the maximum-likelihood
estimator over $\mathcal{H}$ coincides with the minimizer of the empirical
logistic loss
\begin{equation}
  \widehat{L}_S(w, b) \;=\; -\frac{1}{n} \sum_{i=1}^{n}
    \Big[ y_i \log h_{w,b}(x_i) + (1 - y_i) \log\!\big(1 - h_{w,b}(x_i)\big) \Big].
\end{equation}
\end{proposition}

\begin{proof}
The conditional likelihood of the sample under $h_{w,b}$ is
$\prod_i h_{w,b}(x_i)^{y_i} (1 - h_{w,b}(x_i))^{1 - y_i}$, the Bernoulli
probability mass of each label. Taking $-\frac{1}{n}\log(\cdot)$ is a strictly
decreasing transformation, so maximizing the likelihood is equivalent to
minimizing $\widehat{L}_S$.
\end{proof}

\begin{remark}[Convexity]
Writing $p_i = \sigma(z_i)$ with $z_i = \ip{w}{x_i} + b$, each summand is
$\log(1 + e^{z_i}) - y_i z_i$, a convex function of $z_i$ composed with an
affine map of $(w, b)$; hence $\widehat{L}_S$ is convex and full-batch
gradient descent converges to a global minimizer for suitable step sizes.
The gradient takes the simple form
$\nabla_w \widehat{L}_S = \frac{1}{n} \sum_i (p_i - y_i)\, x_i$, which is
precisely the update implemented in the library.
\end{remark}

\begin{remark}[What the class cannot do]
Decision boundaries in $\mathcal{H}$ are hyperplanes
$\{x : \ip{w}{x} + b = 0\}$. Data that is not linearly separable in feature
space bounds the achievable accuracy regardless of training; the companion
writeup (paper 02) measures this empirically by varying class separation.
\end{remark}

\end{document}
